%---------------------------------------------------------------------
% UNIT 7 --- LIMITS, CONTINUITY AND DIFFERENTIATION
% Source: Section 7.14 Self Assessment Questions, pp. 167-169.
%---------------------------------------------------------------------
\chapter{Limits and Differentiation}

\srcnote{Source: \emph{Business Mathematics} 5405/1429, \S7.14 \saq{Self Assessment Questions}, pp.~167--169.
Several sub-parts of Q.6 are illegible in the distributed PDF (radicals and
fraction bars collapsed in the text layer); those are named and skipped rather
than guessed at.}

\begin{keyidea}[The three limit situations, and what to do in each]
\begin{description}
  \item[Substitution works] If $f$ is continuous at $a$, then
        $\lim_{x\to a}f(x)=f(a)$. Try this first, every time.
  \item[You get $\tfrac00$] The zero factors cancel. Factorise, or rationalise
        a surd, then substitute again.
  \item[$x\to\infty$] Divide numerator and denominator by the highest power of
        $x$ in the \emph{denominator}, then use $\tfrac{1}{x^{n}}\to 0$.
\end{description}
And the derivative rules the unit tests:
\[
\dd{}{x}x^{n}=nx^{n-1},\quad
(uv)'=u'v+uv',\quad
\left(\frac uv\right)'=\frac{u'v-uv'}{v^{2}},\quad
\dd{}{x}f(g(x))=f'(g(x))\,g'(x).
\]
\end{keyidea}

%---------------------------------------------------------------------
\section{Estimating limits from a table}

\begin{question}{Q.1}
Estimate the following limits by constructing a table.
\textbf{(a)} $\displaystyle\lim_{x\to 2}(x-3)$ \quad
\textbf{(b)} $\displaystyle\lim_{x\to 1}\frac{x}{\ln x -1}$ \quad
\textbf{(c)} $\displaystyle\lim_{x\to -2}\frac{x^{2}-4}{x+2}$
\end{question}

\begin{solution}
A table approaches $a$ from both sides. If the two columns settle on the same
number, that number is the limit.

\medskip
\textbf{(a)} $f(x)=x-3$ is a polynomial, hence continuous everywhere, so
substitution is legitimate: $\lim_{x\to2}(x-3)=2-3=-1$.

\begin{center}
\small\sffamily
\begin{tabular}{@{}C{1.8cm} C{1.8cm} C{1.0cm} C{1.8cm} C{1.8cm}@{}}
\toprule
$x\to 2^{-}$ & $f(x)$ & & $x\to 2^{+}$ & $f(x)$ \\
\midrule
1.9   & $-1.1$   & & 2.1   & $-0.9$ \\
1.99  & $-1.01$  & & 2.01  & $-0.99$ \\
1.999 & $-1.001$ & & 2.001 & $-0.999$ \\
\bottomrule
\end{tabular}
\end{center}
Both columns close on $-1$.

\medskip
\textbf{(b)} At $x=1$, $\ln 1=0$, so the denominator is $0-1=-1$, not zero.
There is no indeterminacy at all and substitution applies:
\[
\lim_{x\to1}\frac{x}{\ln x-1}=\frac{1}{0-1}=-1 .
\]

\begin{center}
\small\sffamily
\begin{tabular}{@{}C{1.8cm} C{2.2cm} C{1.0cm} C{1.8cm} C{2.2cm}@{}}
\toprule
$x\to 1^{-}$ & $f(x)$ & & $x\to 1^{+}$ & $f(x)$ \\
\midrule
0.9   & $-0.8142$ & & 1.1   & $-1.2159$ \\
0.99  & $-0.9802$ & & 1.01  & $-1.0202$ \\
0.999 & $-0.9980$ & & 1.001 & $-1.0020$ \\
\bottomrule
\end{tabular}
\end{center}

\medskip
\textbf{(c)} Substituting $x=-2$ gives $\tfrac{4-4}{0}=\tfrac00$ ---
indeterminate. Factorise and cancel:
\[
\frac{x^{2}-4}{x+2}=\frac{(x-2)(x+2)}{x+2}=x-2 \quad (x\neq -2),
\]
\[
\lim_{x\to-2}\frac{x^{2}-4}{x+2}=\lim_{x\to-2}(x-2)=-4 .
\]

\begin{center}
\small\sffamily
\begin{tabular}{@{}C{1.8cm} C{1.8cm} C{1.0cm} C{1.8cm} C{1.8cm}@{}}
\toprule
$x\to -2^{-}$ & $f(x)$ & & $x\to -2^{+}$ & $f(x)$ \\
\midrule
$-2.1$   & $-4.1$   & & $-1.9$   & $-3.9$ \\
$-2.01$  & $-4.01$  & & $-1.99$  & $-3.99$ \\
$-2.001$ & $-4.001$ & & $-1.999$ & $-3.999$ \\
\bottomrule
\end{tabular}
\end{center}

\ans{(a) $-1$ \quad (b) $-1$ \quad (c) $-4$}
\end{solution}

\begin{pitfall}
In (c) the function is \emph{undefined} at $x=-2$ --- the cancellation is valid
only for $x\neq-2$. A limit describes the approach, never the value at the
point, and the two differ here.
\end{pitfall}

%---------------------------------------------------------------------
\section{Evaluating limits algebraically}

\begin{question}{Q.2}
Find the indicated limits.
\textbf{(a)} $\displaystyle\lim_{x\to 3}\frac{x^{4}-81}{x^{2}-x-6}$ \quad
\textbf{(b)} $\displaystyle\lim_{x\to \infty}\frac{4-3x^{3}}{x^{3}-1}$ \quad
\textbf{(c)} $\displaystyle\lim_{x\to -7^{-}}\frac{x^{2}+1}{x^{2}-49}$

\smallskip
\textbf{(e)} Find the value of the constant $k$ such that
$\displaystyle\lim_{x\to 2}f(x)$ exists, if
$f(x)=\begin{cases} 2-x, & x<2\\ x^{3}+k(x+1), & x\ge 2\end{cases}$

\smallskip
\textbf{(f)} Use the graph of $f$ to find each of the following limits if they
exist: (i) $\lim_{x\to-1}f(x)$, (ii) $\lim_{x\to 0}f(x)$,
(iii) $\lim_{x\to 1}f(x)$, where
$f(x)=\begin{cases} x^{2}+1, & x<-1\\ 3x+2, & -1\le x\le 1\\ 2-x, & x>1\end{cases}$
\end{question}

\srcnote{Part (d) is illegible in the distributed PDF and is omitted.}

\begin{solution}
\textbf{(a)} Substitution gives $\tfrac{81-81}{9-3-6}=\tfrac00$. Factorise both
parts, using the difference of two squares twice on top:
\[
x^{4}-81=(x^{2}-9)(x^{2}+9)=(x-3)(x+3)(x^{2}+9),
\qquad
x^{2}-x-6=(x-3)(x+2).
\]
The common factor $(x-3)$ cancels:
\[
\lim_{x\to3}\frac{(x-3)(x+3)(x^{2}+9)}{(x-3)(x+2)}
=\lim_{x\to3}\frac{(x+3)(x^{2}+9)}{x+2}
=\frac{(6)(18)}{5}=\frac{108}{5}=21.6 .
\]

\medskip
\textbf{(b)} As $x\to\infty$ this is $\tfrac{-\infty}{\infty}$. Divide top and
bottom by $x^{3}$:
\[
\lim_{x\to\infty}\frac{4-3x^{3}}{x^{3}-1}
=\lim_{x\to\infty}\frac{\tfrac{4}{x^{3}}-3}{1-\tfrac{1}{x^{3}}}
=\frac{0-3}{1-0}=-3 .
\]
Equivalently: for a ratio of polynomials of equal degree, the limit at infinity
is the ratio of leading coefficients, $\tfrac{-3}{1}$.

\medskip
\textbf{(c)} A one-sided limit, approaching $-7$ from the \emph{left}
($x<-7$). The numerator tends to $49+1=50$, a positive number. For the
denominator, $x<-7$ means $|x|>7$, so $x^{2}>49$ and
\[
x^{2}-49\to 0^{+} .
\]
A fixed positive number divided by something positive and shrinking to zero
grows without bound:
\[
\lim_{x\to-7^{-}}\frac{x^{2}+1}{x^{2}-49}=+\infty .
\]
\vcheck{At $x=-7.01$: $x^{2}-49=0.1401$ and $x^{2}+1=50.14$, giving $357.9$.
At $x=-7.001$: the quotient is $3572$. Growing, and positive.}

\medskip
\textbf{(e)} A limit exists at $x=2$ exactly when the two one-sided limits
agree.
\[
\lim_{x\to2^{-}}f(x)=2-2=0,
\qquad
\lim_{x\to2^{+}}f(x)=2^{3}+k(2+1)=8+3k .
\]
Setting them equal:
\[
8+3k=0 \;\Longrightarrow\; k=-\frac83 .
\]
\vcheck{With $k=-\tfrac83$: $f(2^{+})=8-\tfrac83(3)=8-8=0=f(2^{-})$. The limit
exists and equals 0 --- and since $f(2)=0$ too, $f$ is in fact continuous
there.}

\medskip
\textbf{(f)} Check the two one-sided limits at each junction of the piecewise
definition.

\begin{center}
\small\sffamily
\begin{tabular}{@{}c C{3.0cm} C{3.0cm} L{5.0cm}@{}}
\toprule
 & Left-hand limit & Right-hand limit & Conclusion \\
\midrule
(i) $x\to-1$ & $(-1)^{2}+1=2$ & $3(-1)+2=-1$ & unequal $\Rightarrow$ limit does not exist \\
(ii) $x\to 0$ & $3(0)+2=2$ & $3(0)+2=2$ & both from $3x+2$ $\Rightarrow$ limit $=2$ \\
(iii) $x\to 1$ & $3(1)+2=5$ & $2-1=1$ & unequal $\Rightarrow$ limit does not exist \\
\bottomrule
\end{tabular}
\end{center}

At $x=0$ the point lies strictly inside the middle piece, so there is no
junction and the limit is just the value: 2.

\begin{figure}[H]
\centering
\begin{tikzpicture}
\begin{axis}[
  width=11cm, height=6.2cm,
  axis lines=middle, xlabel={$x$}, ylabel={$f(x)$},
  xmin=-3.2, xmax=3.2, ymin=-2.5, ymax=6,
  xtick={-3,-2,...,3}, ytick={-2,0,...,6},
  grid=both, grid style={line width=0.2pt, draw=bmgrey!25},
  tick label style={font=\footnotesize}, label style={font=\small},
  every axis plot/.append style={very thick}
]
\addplot[bmblue,   domain=-3:-1]  {x^2 + 1};
\addplot[bmaccent, domain=-1:1]   {3*x + 2};
\addplot[bmgreen,  domain=1:3]    {2 - x};
% open / closed endpoints marking the two jumps
\addplot[only marks, mark=o, mark size=2.2pt, bmblue,  line width=0.8pt] coordinates {(-1,2)};
\addplot[only marks, mark=*, mark size=2.2pt, bmaccent] coordinates {(-1,-1)};
\addplot[only marks, mark=*, mark size=2.2pt, bmaccent] coordinates {(1,5)};
\addplot[only marks, mark=o, mark size=2.2pt, bmgreen, line width=0.8pt] coordinates {(1,1)};
\node[font=\scriptsize\sffamily, bmred, anchor=west] at (axis cs:-2.9,4.6) {jump at $x=-1$};
\node[font=\scriptsize\sffamily, bmred, anchor=west] at (axis cs:1.35,4.2) {jump at $x=1$};
\end{axis}
\end{tikzpicture}
\caption{The piecewise function of Q.2(f). The limit fails to exist at
$x=-1$ and $x=1$ because the graph jumps; at $x=0$ it is continuous.}
\end{figure}

\ans{(a) $\tfrac{108}{5}=21.6$ \quad (b) $-3$ \quad (c) $+\infty$ \quad
(e) $k=-\tfrac83$ \quad (f) (i) does not exist, (ii) $2$, (iii) does not exist.}
\end{solution}

\begin{pitfall}
In (c) the side matters. From the \emph{right} ($-7<x<7$) we would have
$x^{2}<49$, so $x^{2}-49\to0^{-}$ and the limit would be $-\infty$. A one-sided
limit question is asking you to determine a sign, and that is where the mark is.
\end{pitfall}

%---------------------------------------------------------------------
\section{A population in the long run}

\begin{question}{Q.3}
The population of a certain village $t$ years from now is expected to be
\[
N=15000+\frac{5000}{(t+3)^{2}} .
\]
Determine $\displaystyle\lim_{t\to\infty}N$, i.e.\ the population in the long
run.
\end{question}

\begin{solution}
Only the second term depends on $t$. As $t$ grows, $(t+3)^{2}$ grows without
bound, so the fraction --- a fixed numerator over an exploding denominator ---
shrinks to zero:
\[
\lim_{t\to\infty}\frac{5000}{(t+3)^{2}}=0 .
\]
\[
\lim_{t\to\infty}N=15000+0=15000 .
\]

\medskip
\textbf{Interpretation.} The population declines towards a floor of 15{,}000 and
never falls below it. Today ($t=0$) it is
$15000+\tfrac{5000}{9}\approx 15{,}556$; after 10 years,
$15000+\tfrac{5000}{169}\approx 15{,}030$. The line $N=15000$ is a horizontal
asymptote approached from above.

\ans{$\lim_{t\to\infty}N=\mathbf{15{,}000}$. The population settles at 15{,}000
in the long run, approaching that level from above.}
\end{solution}

%---------------------------------------------------------------------
\section{Average cost per item}

\begin{question}{Q.4}
The cost of manufacturing a certain item is $C(x)=10000+6x$, where $x$ is the
number of items produced. The average cost per item, $\bar{C}(x)$, is obtained
by dividing $C(x)$ by $x$. Determine
\textbf{(a)} $\bar{C}(1500)$, \textbf{(b)} $\bar{C}(10{,}000)$,
\textbf{(c)} $\displaystyle\lim_{x\to 5000}\bar{C}(x)$.
\end{question}

\begin{solution}
\[
\bar{C}(x)=\frac{C(x)}{x}=\frac{10000+6x}{x}=\frac{10000}{x}+6 .
\]
Written this way the structure is plain: a fixed cost of \$10{,}000 spread over
$x$ items, plus \$6 of variable cost that never spreads.

\medskip
\textbf{(a)} $\bar{C}(1500)=\dfrac{10000}{1500}+6=6.667+6=\$12.67$.

\textbf{(b)} $\bar{C}(10000)=\dfrac{10000}{10000}+6=1+6=\$7.00$.

\textbf{(c)} $\bar{C}$ is continuous at $x=5000$ (the only discontinuity is at
$x=0$), so the limit is the value:
\[
\lim_{x\to5000}\bar{C}(x)=\frac{10000}{5000}+6=2+6=\$8.00 .
\]

\medskip
\textbf{Worth adding.} As $x\to\infty$, $\bar{C}(x)\to 6$: average cost falls
towards the variable cost per unit but never reaches it. That is the whole
economics of spreading a fixed cost, and it is a standard follow-up part.

\ans{(a) \$12.67 \quad (b) \$7.00 \quad (c) \$8.00; and
$\bar{C}(x)\to\$6$ as $x\to\infty$.}
\end{solution}

%---------------------------------------------------------------------
\section{Average versus instantaneous rate of change}

\begin{question}{Q.5}
The cost of manufacturing $x$ units of a certain item is
$C(x)=x^{2}-5x+80$.
\textbf{(a)} Find the average rate of change with respect to the number of
items produced between 10 and 15.
\textbf{(b)} Find the instantaneous rate of change when 12 items are produced.
\end{question}

\begin{solution}
\textbf{(a) Average rate of change} is the slope of the chord joining the two
points --- total change divided by total change in $x$:
\[
C(15)=225-75+80=230,\qquad C(10)=100-50+80=130 .
\]
\[
\frac{\Delta C}{\Delta x}=\frac{C(15)-C(10)}{15-10}=\frac{230-130}{5}
=\frac{100}{5}=20 .
\]
Cost rises by \$20 per additional unit, \emph{averaged} over that range.

\medskip
\textbf{(b) Instantaneous rate of change} is the derivative --- the slope of
the tangent at a single point. This is the \emph{marginal cost}:
\[
C'(x)=2x-5, \qquad C'(12)=24-5=19 .
\]
The 13th unit costs about \$19 to produce.

\begin{figure}[H]
\centering
\begin{tikzpicture}
\begin{axis}[
  width=10.5cm, height=6.2cm,
  axis lines=left, xlabel={Units $x$}, ylabel={Cost $C(x)$},
  xmin=6, xmax=18, ymin=60, ymax=280,
  grid=both, grid style={line width=0.2pt, draw=bmgrey!25},
  legend style={font=\footnotesize\sffamily, at={(0.02,0.98)}, anchor=north west,
                draw=bmgrey!50},
  tick label style={font=\footnotesize}, label style={font=\small},
  every axis plot/.append style={very thick}
]
\addplot[bmblue, domain=6:18] {x^2 - 5*x + 80}; \addlegendentry{$C(x)=x^2-5x+80$}
\addplot[bmaccent, domain=9:16, dashed] {130 + 20*(x-10)};
  \addlegendentry{chord, slope $=20$}
\addplot[bmgreen, domain=9.5:14.5, dashed] {164 + 19*(x-12)};
  \addlegendentry{tangent at $x=12$, slope $=19$}
\addplot[only marks, mark=*, mark size=2.2pt, bmred]
  coordinates {(10,130) (15,230) (12,164)};
\end{axis}
\end{tikzpicture}
\caption{The chord from $x=10$ to $x=15$ has slope 20 (average rate); the
tangent at $x=12$ has slope 19 (instantaneous rate).}
\end{figure}

\ans{(a) $20$ per unit \quad (b) $C'(12)=19$ per unit}
\end{solution}

\begin{pitfall}
``Average rate of change'' needs two points and a subtraction; ``instantaneous
rate of change'' needs one point and a derivative. They are different
quantities and give different numbers (20 versus 19 here). Differentiating in
part (a), or substituting into $C$ in part (b), answers the wrong question.
\end{pitfall}

%---------------------------------------------------------------------
\section{Differentiation}

\begin{question}{Q.6 (c)}
Find the derivative of $f(x)=\sqrt{x^{3}+2x^{2}+3}$.
\end{question}

\srcnote{Only sub-part (c) of Q.6 survives legibly in the distributed PDF.
Sub-parts (a), (b), (d), (e) and (f) have collapsed radicals and fraction bars
in the text layer and cannot be transcribed with confidence, so they are
omitted. The drill below exercises the same four rules on clean functions.}

\begin{solution}
This is a chain rule problem: an outer square root wrapped around an inner
polynomial. Write $f(x)=u^{1/2}$ with $u=x^{3}+2x^{2}+3$.
\[
\dd{f}{x}=\frac{1}{2}u^{-1/2}\cdot\dd{u}{x}
=\frac{1}{2\sqrt{u}}\left(3x^{2}+4x\right),
\]
\[
f'(x)=\frac{3x^{2}+4x}{2\sqrt{x^{3}+2x^{2}+3}} .
\]

\ans{$f'(x)=\dfrac{3x^{2}+4x}{2\sqrt{x^{3}+2x^{2}+3}}$}
\end{solution}

\subsection*{Drill on the same four rules}

The following are worked in the style of Q.6 and cover exactly the rules the
unit examines. They are set here, not taken from the book.

\begin{solution}[Power and sum rules]
\[
y=4x^{5}-\frac{7}{x^{2}}+3\sqrt{x}-9
=4x^{5}-7x^{-2}+3x^{1/2}-9 .
\]
Rewriting every term as a power of $x$ \emph{before} differentiating is the
whole technique:
\[
\dd{y}{x}=20x^{4}+14x^{-3}+\frac32 x^{-1/2}
=20x^{4}+\frac{14}{x^{3}}+\frac{3}{2\sqrt{x}} .
\]
\end{solution}

\begin{solution}[Product rule]
\[
y=(5x-2)(3x^{2}+2).
\]
With $u=5x-2$, $u'=5$ and $v=3x^{2}+2$, $v'=6x$:
\[
\dd{y}{x}=u'v+uv'=5(3x^{2}+2)+(5x-2)(6x)
=15x^{2}+10+30x^{2}-12x=45x^{2}-12x+10 .
\]
\vcheck{Expanding first, $y=15x^{3}-6x^{2}+10x-4$, so
$y'=45x^{2}-12x+10$. The two routes agree.}
\end{solution}

\begin{solution}[Quotient rule]
\[
y=\frac{8x^{2}-5}{x^{2}-3}.
\]
With $u=8x^{2}-5$, $u'=16x$ and $v=x^{2}-3$, $v'=2x$:
\[
\dd{y}{x}=\frac{u'v-uv'}{v^{2}}
=\frac{16x(x^{2}-3)-(8x^{2}-5)(2x)}{(x^{2}-3)^{2}}
\]
\[
=\frac{16x^{3}-48x-16x^{3}+10x}{(x^{2}-3)^{2}}
=\frac{-38x}{(x^{2}-3)^{2}} .
\]
\end{solution}

\begin{solution}[Chain rule inside a product]
\[
y=(6x^{3}-x^{2}+3x+5)^{2}.
\]
Outer function squared, inner polynomial differentiated:
\[
\dd{y}{x}=2(6x^{3}-x^{2}+3x+5)^{1}\cdot(18x^{2}-2x+3)
=2(6x^{3}-x^{2}+3x+5)(18x^{2}-2x+3).
\]
\end{solution}

\begin{pitfall}
The quotient rule's numerator is $u'v-uv'$, and the order matters --- swapping
the two terms flips the sign of the whole derivative. A useful memory hook:
the derivative of the \emph{top} comes first.
\end{pitfall}
